\pagebreak
\section{Pregunta generadora}

En la b\'usqueda de informaci\'on sobre nuestra problem\'atica, se propuso la creaci\'on
de un software que se dedique a otorgarle a un usuario la mejor acci\'on que puede
realizar para prevenir fallos de su computadora para que pueda realizar las
acciones que necesita en su d\'{\i}a a d\'{\i}a. Pero en esas circunstancias, se plante\'o
una pregunta que resulta ser el com\'un denominador de los campos que involucra el
software: ¿Qu\'e elementos hay que tener en cuenta para desarrollar el software?

Esta pregunta lo que lleva a un desarrollador es a indagar sobre los conceptos
en lo que es sumamente fatuo para posteriormente, tener un esquema lo
suficientemente satisfactorio en el sentido de que se pueda dar a conocer las
respuestas ante las inc\'ognitas que se genera: Como el qu\'e hacer para crear un
sistema inteligente sea capaz de distinguir entre acciones bajo la condici\'on del
usuario y las de agentes externos, o c\'omo comparar rendimientos en diferentes
periodos de acuerdo al uso y qu\'e consejos puede dar al usuario de tal modo que
no se encuentre con problemas e inc\'ognitas a la hora de solucionar tal
situaci\'on.

Y una de las preguntas que formul\'o \textcite{Trejos2023} para conocer cu\'al es la
mejor manera de aprender a programar en entornos acad\'emicos, como lo son el de
una universidad p\'ublica: ``¿Qu\'e favorece m\'as en el aprendizaje de la
programaci\'on: los fundamentos matem\'aticos y l\'ogicos o los lenguajes de
programaci\'on?'' De esta pregunta nos lleva a pensar en el contexto de realizar
un mantenimiento preventivo en el hardware y software como el preguntarnos sobre
cu\'ales son los medios indicados para dar los resultados que se espera hacia los
usuarios. Y sobre todo, para hacerlos lo suficientemente r\'apidos como para
satisfacer las necesidades actuales.
